\section{Kisah}

Objek tiga dimensi dalam dunia komputasi grafis umumnya direpresentasikan
sebagai \textit{polygon mesh}, yaitu kumpulan segitiga yang membentuk permukaan
suatu objek. Format ini efisien untuk \textit{rendering}, namun kurang praktis
untuk keperluan seperti simulasi fisika, analisis volumetrik, dan manufaktur
berbasis \textit{3D printing}. Salah satu solusi yang umum digunakan adalah
mengonversi representasi \textit{mesh} tersebut menjadi representasi berbasis
\textbf{voksel}, yaitu kubus-kubus kecil berukuran seragam yang menyusun
permukaan objek. Konsep voksel bersifat analog dengan piksel pada gambar dua
dimensi, namun beroperasi dalam ruang tiga dimensi.

Pada Tugas Kecil 2 Strategi Algoritma ini, mahasiswa diminta untuk
mengimplementasikan proses \textit{voxelization} permukaan objek tiga dimensi
menggunakan struktur data \textit{Octree} dan paradigma \textit{Divide and Conquer}.
Program yang dikembangkan diberi nama \textbf{OctoVox}: \textit{Konversi Model 3D
Menjadi Voxel Menggunakan Octree Berbasis Divide and Conquer}.

\section{Latar Belakang}

\textit{Voxelization} adalah proses konversi objek tiga dimensi dari representasi
\textit{polygon mesh} menjadi representasi berbasis voksel. Tantangan utama dalam
proses ini adalah efisiensi komputasi, mengingat ruang tiga dimensi yang harus
diperiksa sangat besar sehingga diperlukan strategi yang cerdas untuk menghindari
pemeriksaan yang tidak perlu.

Pendekatan yang digunakan dalam OctoVox adalah struktur data \textit{Octree},
yaitu sebuah pohon hierarkis di mana setiap \textit{node} merepresentasikan
sebuah kubus dalam ruang tiga dimensi dan dapat memiliki tepat delapan
\textit{node} anak. Dengan struktur ini, ruang tiga dimensi dibagi secara rekursif
hingga resolusi yang diinginkan. Node yang terbukti tidak bersinggungan dengan
permukaan objek akan dipangkas (\textit{pruning}), sehingga pendekatan ini jauh
lebih efisien dibandingkan pendekatan \textit{brute force} yang memeriksa setiap
unit ruang satu per satu.

Proses rekursi tersebut secara langsung mengimplementasikan paradigma
\textit{Divide and Conquer}, di mana masalah besar berupa pemeriksaan seluruh
ruang tiga dimensi dipecah menjadi delapan sub-masalah yang lebih kecil secara
rekursif hingga mencapai resolusi voksel yang ditentukan pengguna.

\section{Deskripsi Tugas}

Program OctoVox menerima sebuah file model tiga dimensi berformat \texttt{.obj}
beserta parameter kedalaman maksimum \textit{Octree}, kemudian menghasilkan file
\texttt{.obj} baru yang merupakan representasi voksel dari permukaan objek masukan.
Program dikembangkan dalam bahasa pemrograman Go dan dapat dijalankan pada sistem
operasi Windows maupun Linux. Program ini dikerjakan secara berkelompok oleh dua
orang mahasiswa.

\subsection{Masukan}

Program menerima dua buah parameter masukan melalui antarmuka CLI:

\begin{enumerate}
    \item \textbf{Path file \texttt{.obj}}: jalur menuju file model tiga dimensi
    yang ingin dikonversi. Format \texttt{.obj} yang valid hanya memuat dua jenis
    baris, yaitu:
    \begin{itemize}
        \item \texttt{v x y z}, yang mendefinisikan sebuah titik (\textit{vertex})
        pada koordinat $(x, y, z)$ dalam ruang tiga dimensi.
        \item \texttt{f i j k}, yang mendefinisikan sebuah segitiga (\textit{face})
        dengan menghubungkan \textit{vertex} ke-$i$, ke-$j$, dan ke-$k$.
    \end{itemize}
    Berikut adalah contoh konten file \texttt{.obj} yang valid:
    \begin{verbatim}
v 2.229345 -0.992723 -0.862826
v 2.292449 -0.871852 -0.882400
v 2.410367 -0.777999 -0.841105
f 1 2 3
    \end{verbatim}
    Program akan mendeteksi dan menolak file yang tidak memenuhi format tersebut
    dengan menampilkan pesan kesalahan yang informatif kepada pengguna.

    \item \textbf{Kedalaman maksimum Octree ($N$)}: bilangan bulat positif yang
    menentukan resolusi voksel pada keluaran. Semakin besar nilai $N$, semakin
    kecil ukuran voksel yang dihasilkan dan semakin tinggi tingkat detail objek
    yang direpresentasikan. Ukuran sisi voksel memenuhi persamaan berikut:
    \begin{equation}
        s_{\text{voksel}} = \frac{s_{\text{bounding box}}}{2^N}
    \end{equation}
\end{enumerate}

\subsection{Keluaran}

Program menghasilkan dua bentuk keluaran secara bersamaan setelah proses
\textit{voxelization} selesai dijalankan:

\begin{enumerate}
    \item \textbf{File \texttt{.obj} hasil konversi}: file model tiga dimensi
    baru yang tersusun dari voksel-voksel berukuran seragam. Setiap voksel
    direpresentasikan sebagai sebuah kubus dengan \textbf{8 \textit{vertex}} dan
    \textbf{12 \textit{face}}, yaitu dua segitiga per sisi dikalikan enam sisi
    kubus. Apabila terdapat $V$ voksel dalam hasil konversi, maka file keluaran
    memuat $8V$ baris \textit{vertex} dan $12V$ baris \textit{face}.

    \item \textbf{Laporan pada CLI}: informasi statistik yang ditampilkan segera
    setelah proses selesai, memuat hal-hal berikut:
    \begin{itemize}
        \item Banyaknya voksel yang terbentuk.
        \item Banyaknya \textit{vertex} dan \textit{face} pada file keluaran.
        \item Statistik jumlah \textit{node} Octree yang terbentuk per kedalaman:
        \begin{verbatim}
1 : <banyaknya node dengan depth 1>
2 : <banyaknya node dengan depth 2>
...
N : <banyaknya node dengan depth N>
        \end{verbatim}
        \item Statistik jumlah \textit{node} yang tidak perlu ditelusuri
        (\textit{pruned}) per kedalaman, dengan format yang sama.
        \item Kedalaman \textit{Octree} yang digunakan.
        \item Lama waktu eksekusi program.
        \item \textit{Path} penyimpanan file \texttt{.obj} keluaran.
    \end{itemize}
\end{enumerate}
